% Generated draft: requires mathematical and source review.
\documentclass[11pt]{article}
\usepackage[margin=25mm]{geometry}
\usepackage{fontspec}
\setmainfont{Latin Modern Roman}
\usepackage{amsmath,amssymb,mathtools,unicode-math}
\setmathfont{Latin Modern Math}
\usepackage{enumitem,needspace,xurl,hyperref}
\setlength{\parindent}{0pt}
\setlength{\parskip}{5pt}
\setlength{\emergencystretch}{3em}
\newcommand{\N}{\mathbb{N}}
\newcommand{\Z}{\mathbb{Z}}
\newcommand{\Q}{\mathbb{Q}}
\newcommand{\R}{\mathbb{R}}
\newcommand{\C}{\mathbb{C}}
\newcommand{\F}{\mathbb{F}}
\newcommand{\A}{\mathbb{A}}
\newcommand{\PP}{\mathbb{P}}
\newcommand{\E}{\mathbb{E}}
\newcommand{\eps}{\varepsilon}
\newcommand{\dd}{\,\mathrm{d}}
\begin{document}
\section*{UnsolvedMath ID: 20002015}
\textbf{Source problem code:} catalogue entry
\subsection{Definitions and mathematical statement}
Problem 2 asks for natural tensor-densities with the affine conformal transformation law labeled (43) in the source. That preceding formula is not reproduced. The paired aid uses the following convention. On a conformal \(n\)-manifold, let \(\widehat g=e^{2\omega}g\), let \(\mathcal E[w]\) denote weight-\(w\) densities, and let \(\mathcal T[-1]=\mathcal T\otimes\mathcal E[-1]\) be the weighted standard tractor bundle. The scale-associated tractor \(I_g\) and associated operator \(D\) are normalized so their scale-change law gives the affine construction below; their complete formulas must be checked against (43).

Let \(P:\mathcal T[-1]\to\mathcal E[-n]\) be a natural conformally invariant linear differential operator. The aid reports
\[
 A_g=PI_g,\qquad L=PD,\qquad
 A_{\widehat g}=A_g+L\omega,\qquad L1=0,
\]
as an identity of density sections. In scalar representatives this reads
\[
 e^{n\omega}a_{\widehat g}=a_g+L_g\omega.
\]
The operator \(L\) is conformally invariant in the corresponding density sense. This is a construction of examples, not an exhaustive classification.

For even dimension in a setting with critical \(Q\)-curvature and GJMS operator \(P_n\), use the convention
\(e^{n\omega}q_{\widehat g}=q_g+(P_n)_g\omega\). If \(L\) is formally self-adjoint on a closed manifold, then for any constant \(\alpha\),
\[
 \widetilde q_g=q_g+\alpha a_g,\qquad
 \widetilde P_g=(P_n)_g+\alpha L_g
\]
satisfy the same affine law, and the integral of \(\widetilde q_g\) is conformally invariant. A general correction need not have zero integral; its integral is invariant. The aid reports the quadratic conformal cocycle
\[
 S_g(\omega)=\int_M\left(
 \omega\widetilde q_g+\tfrac12\omega\widetilde P_g\omega
 \right)dv_g,
\]
\[
 S_g(\omega+\eta)=S_g(\omega)+S_{e^{2\omega}g}(\eta).
\]
Self-adjointness, annihilation of constants and the boundary-free integration convention are essential here.

The workshop example is \(P=\iota(D)|C|^2\), where \(C\) is Weyl curvature, \(|C|^2\) acts by multiplication and \(\iota(D)\) denotes tractor contraction. The source discusses this formal construction in any dimension and singles out dimension six for the integrable weight \(-n\). Its six-dimensional scalar representatives, in the workshop sign and normalization convention, are
\[
 a_g=4\Delta_g|C_g|^2,\qquad
 L_g=-16\delta_g\bigl(|C_g|^2d\bigr),
\]
\[
 e^{6\omega}a_{\widehat g}
 =4\Delta_g|C_g|^2-16\delta_g\bigl(|C_g|^2d\omega\bigr).
\]
The extracted trailing ``21'' in the first formula is treated as a page artifact, as the supplied summary reconstructs the expression. The divergence-form operator is formally self-adjoint. On every closed six-manifold,
\(\int_M a_g\,dv_g=0\), so this example leaves the total \(Q\) value unchanged while changing the local prescription operator wherever Weyl curvature is nonzero. It relates \(Q_6+\alpha a_g\) to a modification of \(P_6\), rather than asserting all the original \(Q_6\) properties are identical.

The aid also names later Branson--Gover \(Q_k\)/Pontrjagin constructions as further scalar families, but supplies no formulas or exhaustive classification. The original record ends mid-sentence, so no missing final assertion is reconstructed.
\subsection{Short English statement}
The supplied tractor construction produces densities with an affine conformal law and a conformally invariant correction operator. Under self-adjointness it yields modified \(Q\)-curvature and a quadratic cocycle. The six-dimensional Weyl-square example changes the local operator but has zero total correction on closed manifolds. The requested general classification is not supplied.
\subsection{Sources}
\begin{itemize}
\item Local source record: catalogue entry, catalogue entry; canonical \texttt{aim-geometry-notes.json}, \texttt{notes[352]}, Problem 2. Title: ``Affine tractor factories and Q-like tensor-densities.''
\item AIM notes \emph{Conformal Structure in Geometry, Analysis, and Physics}, version 15 October 2003, section ``Other routes to \(Q\) and its variants,'' supplied reference \url{https://aimath.org/WWN/confstruct/confstruct.pdf}. The supplied verification refers to printed pages 20--21; this draft does not access a PDF.
\item Source constructions: weighted tractors or the ambient metric, Weyl-square contraction and the GJMS operator. Supplied research summary also names Branson--Gover \(Q_k\)/Pontrjagin constructions without full bibliographic entries. Attempt 1 gives the self-adjoint cocycle synthesis and zero-total six-dimensional observation.
\item Registry dated 14 August 2026: ``partially\_solved,'' ``PARTIAL-PROGRESS,'' statement status ``exact''; attempt partial, rigor ``proved,'' novelty confidence low. The machine-generated, non-peer-reviewed aid requires independent verification before citation; the canonical extraction is truncated.
\end{itemize}
\end{document}
